给定有限样本，我们想在未观测位置重建函数。插值要求曲线逐点命中数据，逼近则允许偏离样本以换取整体误差更小或模型更简单。

取 \(f(x)=1/(1+25x^2)\)，在 \([-1,1]\) 上均匀取 11 个点插值。10 次插值多项式在两端的振荡幅度超过 50——而在区间的中部，误差可控制在 \(10^{-4}\) 以内。同样这些点，若改用 Chebyshev 节点（不等距，两端密、中间疏），10 次多项式的最大误差就降到 \(10^{-2}\) 量级以下。节点在哪里放着，比增加次数在区间中部扩大精度重要得多。

高次多项式虽然自由度多，却可能因节点几何在端点剧烈振荡；样条通过分段低次结构获得局部性。这个单元关注数据、函数空间、误差范数与节点选择怎样共同决定结果。

\begin{enumerate}
\tightlist
\item
  一个插值多项式的两种构造
\item
  Runge 现象与节点选择
\item
  从精确插值到整体逼近
\item
  三次样条：用局部性换稳定
\end{enumerate}

第一次读先把三件事分开：插值要求穿过给定点，逼近关心整个区间的误差，稳定求值则关心计算过程会不会放大舍入和数据扰动。Lagrange、Newton 和样条不是三份要背的公式，而是在回答不同问题：节点能否选择、数据会不会增加、改一个点应影响多远。

前一单元：\hyperref[ch:062]{``非线性方程求根：可靠性、速度与信息成本''}。

\subsection{本章篇目}

以下篇目按概念依赖排列；若从具体问题进入，也可以先打开最接近当前困惑的一篇，再沿篇内链接回补。

\begin{itemize}
\tightlist
\item \hyperref[sec:08-Numerical-Analysis/03-Interpolation-and-Approximation/01]{``一个插值多项式的两种构造''}；
\item \hyperref[sec:08-Numerical-Analysis/03-Interpolation-and-Approximation/02]{``Runge 现象与节点选择''}；
\item \hyperref[sec:08-Numerical-Analysis/03-Interpolation-and-Approximation/03]{``从精确插值到整体逼近''}；
\item \hyperref[sec:08-Numerical-Analysis/03-Interpolation-and-Approximation/04]{``三次样条：用局部性换稳定''}。
\end{itemize}
